\section{The Supergraph}
\label{sec:ch06-the-supergraph}
\index{supergraph|(}%

Start with the picture that everything from here to
chapter~\ref{ch:the-gateway-market-and-the-future} hangs on.

\begin{graphqlsdl}
# -- Catalog subgraph --
type Product @key(fields: "id") {
  id: ID!
  title: String!
  price: Money!
}

# -- Reviews subgraph --
type Product @key(fields: "id") {
  id: ID!
  reviews: [Review!]!
}

# -- The supergraph, after composition --
type Product {
  id: ID!
  title: String!
  price: Money!
  reviews: [Review!]!
}
\end{graphqlsdl}

Two subgraphs, each defining a \code{Product} type with different fields. A
composer reads both schemas, notices they share a \code{@key}, and merges them
into one \code{Product} with four fields. A client that sends a query selecting
\code{title}, \code{price} and \code{reviews} gets an answer from two services
without knowing there were two. The router knew: it fetched the Catalog's fields
first, then used the \code{id} to ask Reviews for the rest.

\index{router!role}%
The router is the single point of contact, and the single point of contact is
the whole trick. The client does not address a subgraph directly. It does not
know which subgraph owns which field. It sends one document to one endpoint, and
the router's query planner produces a plan that says which subgraph answers
which selection, in which order, with which identifiers carried between them.
Chapter~\ref{ch:how-a-federated-query-actually-runs} is that plan in motion;
this chapter is the static structure that makes it possible.

\index{subgraph!definition}%
A subgraph is a normal GraphQL service with two extra responsibilities. It has
to publish its own schema without knowing about the others, and it has to answer
two queries the router will send: \code{\_service}, which returns the SDL the
router needs for composition, and \code{\_entities}, which receives a list of
type-and-key pairs and returns the matching objects. Chapter~\ref{ch:how-a-federated-query-actually-runs}
covers both in detail. For now, all that matters is that a subgraph declares its
own types and fields, marks which ones are shareable, and says which are
entities the router can ask for.

\index{supergraph!schema}%
The supergraph schema is what the client sees, and it is not one subgraph's
schema with additions. It is the output of composition: the union of every type,
every field and every directive from every subgraph, merged according to rules
that sometimes reject the whole arrangement. A field that appears in two
subgraphs without \code{@shareable} on both is a composition error. A field of
\code{Product} that a query can reach but no subgraph can resolve is a
composition error. The composer is the gatekeeper for the claim that the
supergraph works, and section~\ref{sec:ch06-how-composition-thinks} is its
rules.

\index{Mosaic!as federated supergraph}%
Mosaic, from chapter~\ref{ch:the-first-cut-extracting-catalog} onward, is six
subgraphs behind one router. Catalog owns products and their attributes.
Reviews owns reviews. Orders owns orders and their lines. Inventory owns stock.
Pricing owns prices and currency. Accounts owns customers and their profiles.
Each is a service with its own database, its own deployment, and its own slice
of the schema, and together they answer the same queries the monolith did. The
rest of Part~II is about how that arrangement is built, and this chapter is the
language it is described in.

\index{supergraph|)}%
